\section*{अध्याय \ref{ch:b1:replication-mitosis} --- DNA प्रतिकृतिकरण और समसूत्री विभाजन}

\begin{solution}{exo:b1:replication-mitosis:1}
पॉलिमरेज़ \omterm{def:b1:nucleic-acids:nucleotide}{न्यूक्लियोटाइड} केवल किसी $3'$ सिरे पर जोड़ती हैं, इसलिए नई
लड़ियाँ $5' \to 3'$ बढ़ती हैं; और दोनों टेम्पलेट प्रतिसमांतर हैं, इसलिए
किसी द्विशाखिका पर एक नई लड़ी द्विशाखिका की ओर लगातार बढ़ सकती है जबकि
दूसरी को उससे दूर बढ़ना पड़ता है, और ज्यों-ज्यों और टेम्पलेट खुलता है
त्यों-त्यों टुकड़ों में फिर से शुरू होना पड़ता है।
\end{solution}

\begin{solution}{exo:b1:replication-mitosis:2}
हेलिकेज़ (खोलती है), एकलड़ी-बंधनकारी \omterm{def:b1:proteins:peptide}{प्रोटीन} (लड़ियों को अलग थामे
रखते हैं), टोपोआइसोमरेज़ (आगे का ऐंठन ढीला करती है), प्राइमेज़ (\omterm{def:b1:nucleic-acids:chain}{RNA} प्राइमर),
\omterm{def:b1:replication-mitosis:fork}{DNA पॉलिमरेज़} (प्राइमर बढ़ाती और शुद्धि जाँचती है), एक दूसरी पॉलिमरेज़
(प्राइमर बदलती है), लाइगेज़ (टुकड़े सील करती है)।
\end{solution}

\begin{solution}{exo:b1:replication-mitosis:3}
G$_1$: $q$, एक \omterm{def:b1:replication-mitosis:mitosis}{क्रोमैटिड} वाले 2n \omterm{def:b1:genomes:genome}{गुणसूत्र}। G$_2$: $2q$, दो
\omterm{def:b1:replication-mitosis:mitosis}{क्रोमैटिडों} वाले 2n \omterm{def:b1:genomes:genome}{गुणसूत्र}। \omterm{def:b1:replication-mitosis:mitosis}{समसूत्री विभाजन} के बाद: प्रति संतति $q$,
और हर एक में 2n \omterm{def:b1:genomes:genome}{गुणसूत्र}।
\end{solution}

\begin{solution}{exo:b1:replication-mitosis:4}
पूर्वावस्था (संघनन, \omterm{def:b1:replication-mitosis:mitosis}{तर्कु} बनता है), पूर्व-मध्यावस्था (आवरण टूटता है,
गतिबिंब जुड़ते हैं), मध्यावस्था (विषुवत् पर पंक्तिबद्धता), पश्चावस्था
(\omterm{def:b1:replication-mitosis:mitosis}{क्रोमैटिड} अलग होते हैं), अंत्यावस्था (आवरण फिर बनते हैं), फिर
कोशिकाद्रव्य-विभाजन।
\end{solution}

\begin{solution}{exo:b1:replication-mitosis:5}
आधा \omterm{def:b1:genomes:genome}{जीनोम} पश्चगामी-लड़ी का संश्लेषण है: $\num{6.4e9}/150 =
\num{4.3e7}$ टुकड़े; और उतने ही प्राइमर, हटाव तथा जोड़।
\end{solution}

\begin{solution}{exo:b1:replication-mitosis:6}
$\num{6.4e9}\times 10^{-5} = \num{64000}$; $\times 10^{-7}$: 640;
$\times 10^{-9}$: यानी प्रति विभाजन लगभग 6 उत्परिवर्तन।
\end{solution}

\begin{solution}{exo:b1:replication-mitosis:7}
$5000/80 = 62$ विभाजन। किसी अर्बुद को बिना सीमा विभाजित होना पड़ता है;
टीलोमरेज़ के बिना उसके टीलोमर चुक जाते और उसकी \omterm{def:b1:cell-unit-of-life:cell}{कोशिकाएँ} रुक जातीं या मर
जातीं, इसलिए लगभग हर कर्कट उस \omterm{def:b1:enzymes:enzyme}{एंजाइम} को फिर सक्रिय कर देता है।
\end{solution}

\begin{solution}{exo:b1:replication-mitosis:8}
चक्र $= 1\,\mathrm{h}/0.05 = \qty{20}{h}$; S $= 0.30\times 20 =
\qty{6}{h}$।
\end{solution}

\begin{solution}{exo:b1:replication-mitosis:9}
G$_1$: 8 \omterm{def:b1:genomes:genome}{गुणसूत्र}, 8 \omterm{def:b1:replication-mitosis:mitosis}{क्रोमैटिड}, $q$। G$_2$: 8, 16, $2q$। मध्यावस्था:
8, 16, $2q$। संतति: 8, 8, $q$।
\end{solution}

\begin{solution}{exo:b1:replication-mitosis:10}
पिछला चक्कर खत्म होने से पहले ही उद्गम पर हर \qty{20}{min} में नए चक्कर
शुरू हो जाते हैं: प्रतिकृतिकरण अतिव्यापी है, और किसी नवजात \omterm{def:b1:cell-unit-of-life:cell}{कोशिका} के पास
पहले से ही अंशतः प्रतिकृत \omterm{def:b1:nucleic-acids:chain}{DNA} होता है। विभाजन पर जो \omterm{def:b1:genomes:genome}{गुणसूत्र} अपना चक्कर
पूरा कर रहा होता है वह \qty{40}{min} पहले शुरू हुआ था, \qty{20}{min}
पहले शुरू हुए दूसरे चक्कर की द्विशाखिकाएँ आधे रास्ते हैं, और तीसरा अभी
शुरू ही हो रहा है: विभाजित होती \omterm{def:b1:cell-unit-of-life:cell}{कोशिका} में 4 उद्गम और 6 द्विशाखिकाएँ।
\end{solution}

\begin{solution}{exo:b1:replication-mitosis:11}
कोई \omterm{def:b1:replication-mitosis:mitosis}{तर्कु} नहीं बनता: \omterm{def:b1:genomes:genome}{गुणसूत्र} संघनित तो होते हैं पर न पंक्तिबद्ध हो सकते
हैं न अलग, और \omterm{def:b1:cell-unit-of-life:cell}{कोशिका} मध्यावस्था-जैसी अवस्था में तर्कु-जाँचबिंदु पर रुक
जाती है। यदि वह बिना विभाजित हुए \omterm{def:b1:replication-mitosis:cycle}{अंतरावस्था} में फिसल जाए तो उसके पास
4n \omterm{def:b1:genomes:genome}{गुणसूत्र} (चतुर्गुणित) होते हैं। संघनित, अलग \omterm{def:b1:genomes:genome}{गुणसूत्रों} वाली रुकी
\omterm{def:b1:cell-unit-of-life:cell}{कोशिकाएँ} ठीक वही हैं जो किसी \omterm{def:b1:genomes:karyotype}{गुणसूत्र-चित्र} को चाहिए: उन्हें जमा करने के
लिए यह औषधि लगाई जाती है।
\end{solution}

\begin{solution}{exo:b1:replication-mitosis:12}
प्रतिकृतिकरण यह पक्का करता है कि हर \omterm{def:b1:genomes:genome}{गुणसूत्र} दो समरूप \omterm{def:b1:replication-mitosis:mitosis}{क्रोमैटिड} बने,
एक अरब में एक गलती तक; और \omterm{def:b1:replication-mitosis:mitosis}{तर्कु} यह पक्का करता है कि हर संतति को हर जोड़े
में से एक मिले, हर \omterm{def:b1:genomes:genome}{गुणसूत्र} को दोनों ध्रुवों से जोड़कर और पश्चावस्था को
तब तक रोककर जब तक हर एक तनाव में न आ जाए। G$_2$ \omterm{prop:b1:replication-mitosis:checkpoints}{जाँचबिंदु} यह पक्का करता
है कि गिनती शुरू होने से पहले नकल पूरी हो चुके; और तर्कु-जाँचबिंदु यह कि
अलगाव से पहले गिनती बैठ चुके। जो \omterm{def:b1:cell-unit-of-life:cell}{कोशिका} गलत बाँटती है उसके पास एक
\omterm{def:b1:genomes:genome}{गुणसूत्र} अधिक या कम रह जाता है (असुगुणिता): यानी अक्षर-अक्षर सही कोई नकल,
जो गलत पते पर पहुँचाई गई --- और यही ऐसी अधिकांश \omterm{def:b1:cell-unit-of-life:cell}{कोशिकाओं} को मार डालता है
तथा अधिकांश कर्कटों का लक्षण है।
\end{solution}

\begin{solution}{pb:b1:replication-mitosis:1}
\textbf{1.} $50\times 28\,800 = \qty{1.44e6}{bp}$।
\textbf{2.} \qty{2.88e6}{bp}।
\textbf{3.} $\num{6.4e9}/\num{2.88e6} = 2220$ उद्गम।
\textbf{4.} \qty{2.9}{Mb}, यानी उद्गमों के बीच \omterm{def:b1:nucleic-acids:chain}{DNA} की $2.9\times 10^6\times 0.34\,\mathrm{nm} =
\qty{0.98}{mm}$।
\textbf{5.} $\num{6.4e9}/\num{40000} = \qty{160}{kb}$ का अंतराल; हर
द्विशाखिका लगभग \qty{80}{kb} चलती है, यानी आधे घंटे का काम।
\textbf{6.} \qty{50}{bp/s} पर प्रति द्विशाखिका \qty{125}{Mb}: $\num{2.5e6}$ s,
यानी 29 दिन।
\textbf{7.} द्विशाखिका की चाल रसायन से और क्रोमैटिन से सीमित है
(बैक्टीरिया के सापेक्ष बीस गुनी धीमी); उद्गम बिना सीमा बढ़ाए जा सकते हैं,
इसलिए S प्रावस्था द्विशाखिकाओं को दौड़ाकर नहीं, शुरुआती बिंदु जोड़कर
छोटी की जाती है।
\textbf{8.} पश्चगामी संश्लेषण आधे \omterm{def:b1:genomes:genome}{जीनोम} को ढँकता है, यानी $\num{3.2e9}$
\omterm{def:b1:nucleic-acids:nucleotide}{न्यूक्लियोटाइड}: $\num{2.1e7}$ टुकड़े (हर उद्गम की दोनों द्विशाखिकाओं की
अपनी पश्चगामी लड़ी होती है, और दोनों आधे मिलकर एक \omterm{def:b1:genomes:genome}{जीनोम} जितने बनते हैं)।
\textbf{9.} $\num{2.1e7} + 2\times\num{40000} = \num{2.1e7}$ प्राइमर।
\textbf{10.} लगभग $\num{2.1e7}$ जोड़ (प्रति टुकड़ा एक) जोड़
द्विशाखिकाओं की हर भेंट पर एक।
\textbf{11.} $2\times\num{6.4e9} = \num{1.28e10}$ \omterm{def:b1:nucleic-acids:nucleotide}{न्यूक्लियोटाइड}।
\textbf{12.} $\num{2.56e10}$ \omterm{def:b1:water-small-molecules:atp}{ATP}; \qty{28800}{s} भर, प्रति सेकंड
$\num{8.9e5}$।
\textbf{13.} यानी \omterm{def:b1:cell-unit-of-life:cell}{कोशिका} के ATP-आवर्तन का लगभग \qty{0.1}{\%}:
बहुलकन सस्ता है; महँगा है \omterm{def:b1:nucleic-acids:nucleotide}{न्यूक्लियोटाइड} और \omterm{def:b1:genomes:nucleosome}{हिस्टोन} बनाना।
\textbf{14.} \num{64000}, 640 और 6.4 गलतियाँ।
\textbf{15.} $1500\times 10^{-9} = \num{1.5e-6}$: यानी प्रति विभाजन सात
लाख में एक \omterm{def:b1:genomes:gene}{जीन}।
\textbf{16.} $10^{16}\times 6.4 = \num{6.4e16}$ उत्परिवर्तन --- यानी
\omterm{def:b1:genomes:genome}{जीनोम} का हर संभव एकल बदलाव कई-कई बार, जो $10^{13}$ \omterm{def:b1:cell-unit-of-life:cell}{कोशिकाओं} में बिखरा
है: लगभग सारे अकूटलेखी \omterm{def:b1:nucleic-acids:chain}{DNA} में या ऐसी \omterm{def:b1:cell-unit-of-life:cell}{कोशिकाओं} में पड़ते हैं जो झड़
जाती हैं, और कोई उत्परिवर्तित \omterm{def:b1:cell-unit-of-life:cell}{कोशिका} अरबों में से एक है; केवल किसी एक ही
\omterm{def:b1:cell-unit-of-life:cell}{कोशिका} के कुछ ही संयोजन मायने रखते हैं।
\textbf{17.} सौ गुना ($10^{-7}$): कर्कट बनाने वाले उत्परिवर्तनों का जो
क्रम $10^{-9}$ पर जमा होने में दशकों लेता है, वह कहीं जल्दी
जमा हो जाता है --- बेमेल-मरम्मत के वंशागत दोष
जवानी में ही बृहदांत्र-कर्कट कर देते हैं।
\textbf{18.} \qty{1000}{bp/s} पर प्रति द्विशाखिका $\num{2.3e6}$ युग्म:
\qty{2300}{s}, यानी \qty{38}{min}।
\textbf{19.} लगभग दो चक्कर अतिव्यापी; कोई नवजात \omterm{def:b1:cell-unit-of-life:cell}{कोशिका}
दो उद्गम ढोती है (हर एक की एक बार नकल हो चुकी) और उन पर द्विशाखिकाएँ आधे
रास्ते: उसका \omterm{def:b1:genomes:genome}{गुणसूत्र} जन्म पर ही अंशतः प्रतिकृत होता है।
\textbf{20.} प्रति चक्कर $\num{4.6e6}/1500 = 3070$ टुकड़े, यानी प्रति
सेकंड लगभग 1.3.
\textbf{21.} प्रति चक्कर $\num{4.6e6}\times 10^{-9} = \num{4.6e-3}$
गलतियाँ: यानी 200 में लगभग एक संतति कोई नया उत्परिवर्तन ढोती है।
\textbf{22.} प्रति मिलीलीटर $10^9\times\num{4.6e-3} = \num{4.6e6}$ नए
उत्परिवर्तन; और किसी दिए गए \qty{1}{kb} \omterm{def:b1:genomes:gene}{जीन} पर $\num{4.6e6}\times
1000/\num{4.6e6} = 1000$ बार चोट पड़ती है।
\textbf{23.} प्रति मिलीलीटर प्रति विभाजन किसी भी दिए गए \omterm{def:b1:genomes:gene}{जीन} में हज़ार
स्वतंत्र उत्परिवर्तनों के साथ, हर संभव एकल बदलाव --- वे भी जो
प्रतिरोध देते हैं --- प्रतिजैविक लगाए जाने से पहले ही किसी बड़ी
संवर्धि में मौजूद रहता है; औषधि चुनती है, बनाती
नहीं।
\textbf{24.} मनुष्य: $\num{6.4e9}$ युग्म, \num{40000} उद्गम,
\qty{50}{bp/s}, \qty{8}{h}। बैक्टीरियम: $\num{4.6e6}$ युग्म, एक उद्गम,
\qty{1000}{bp/s}, \qty{40}{min}। यानी हज़ार गुना बड़ा \omterm{def:b1:genomes:genome}{जीनोम}, बीस गुना
धीमी द्विशाखिकाओं से नकल किया जाता, और केवल बारह गुने समय में:
यह अंतर समांतर में काम करते दसियों हज़ार उद्गमों का
है।
\textbf{25.} कम से कम लगभग \num{2200} उद्गम, यदि सब एक साथ चलें;
असल में लगभग \num{40000}, हर \qty{160}{kb} पर एक, जो एक के बाद एक
चलते हैं ताकि हर द्विशाखिका लगभग \qty{80}{kb} चले।
\end{solution}
